% Source: Independent derivation; cross-checked with KACTL geometry.
\subsection{直线求交、投影、反射与距离}
\label{geo:line}
\AlgorithmInfo
    {构造实数交点，求点到直线／线段的距离。}
    {$O(1)$。}
    {依赖实数点；Line 两构造点必须不同，方向内部归一化。}
\subsubsection*{原理}

直线记为 $p+tv$，v 为单位方向。两直线求交时，令 $d=v_1\times v_2$；非平行时
\[
 t=\frac{(p_2-p_1)\times v_2}{d}.
\]
平行时再检查 p2 到第一条线的有向距离，区分重合与平行分离。

投影为 $p_0+v((q-p_0)\cdot v)$，反射点是投影的两倍减原点。点到直线距离为 $|v\times(q-p_0)|$。点到线段时投影参数限制在 [0,1]；端点相同则直接求点间距离，不做除法。

\subsubsection*{使用教程}
\InlineCode{Line(a,b)} 接收两个\textbf{不同}的实数点，表示经过 a、沿 a 到 b 方向的无限直线；成员 p 是线上点，v 是单位方向。\InlineCode{l.side(q)} 返回有符号距离，左正右负，0 在直线上；交换构造端点会交换保留侧，但不影响距离和投影。构造后不要单独将 v 改成非单位向量。

\InlineCode{line_intersection(a,b)} 返回结构体：kind=0 为平行异线，1 为唯一交点，2 为重合；\textbf{仅当 kind=1 时读取 p}。它求的是无限直线交点，不保证落在原先的端点线段内；近乎平行时也要关注交点误差。

\InlineCode{projection(q,l)} 返回垂足，\InlineCode{reflection(q,l)} 返回对称点，\InlineCode{point_line_distance(q,l)} 返回到无限直线的非负距离。\InlineCode{point_segment_distance(q,a,b)} 则返回到闭线段的距离，会限制垂足范围，并允许 a=b。均无需额外初始化，不修改传入对象。
\begin{lstlisting}[style=mtycpp,numbers=none]
Line a(Point{0,0},Point{2,2}),b(Point{0,2},Point{2,0});
auto hit=line_intersection(a,b);
if(hit.kind==1){Point intersection=hit.p;}
Point p=projection(Point{2,0},a),q=reflection(Point{2,0},a);
ld side=a.side(Point{2,0});
ld dl=point_line_distance(Point{3,1},a);
ld ds=point_segment_distance(Point{3,1},Point{0,0},Point{2,0});
\end{lstlisting}
\needspace{7\baselineskip}
\subsubsection*{模板代码}
\CodeFile{code/geometry/line.cpp}
